\chapter{二重积分}

\section{原题页索引}
本章原题对应做题本第 96--111 页。第 109 页原题题号从 \(34\) 跳到 \(36\)，原书可见页面中未出现第 \(35\) 题。

\begin{center}
\includegraphics[width=.92\textwidth]{figures/original_pages/page_096.png}
\end{center}

\section{详细解析}

\begin{problemblock}
\textbf{1.} 设 \(f(x,y)\) 连续，判断下列二次积分换序结果。
\begin{enumerate}
\item \(\displaystyle \int_1^2 dx\int_x^2 f(x,y)\,dy+\int_1^2dy\int_y^{4-y}f(x,y)\,dx\).
\item \(\displaystyle \int_{\pi/2}^{\pi}dx\int_{\sin x}^1 f(x,y)\,dy\).
\end{enumerate}
\end{problemblock}

\begin{solutionblock}
\analysis{换序题的核心是还原积分区域。本题两问都是选择题，答案分别为 C、B。}
\begin{enumerate}
\item 第一部分区域为
\[
1\le x\le2,\qquad x\le y\le2,
\]
换成 \(y\) 先固定时是 \(1\le y\le2,\,1\le x\le y\)。第二部分已经是
\[
1\le y\le2,\qquad y\le x\le4-y.
\]
合并后
\[
1\le y\le2,\qquad 1\le x\le4-y.
\]
故
\[
\boxed{\int_1^2dy\int_1^{4-y}f(x,y)\,dx}.
\]
\item 原区域为
\[
\frac{\pi}{2}\le x\le\pi,\qquad \sin x\le y\le1.
\]
当 \(0\le y\le1\) 时，在 \([\pi/2,\pi]\) 上满足 \(\sin x\le y\) 的 \(x\) 为
\[
\pi-\arcsin y\le x\le\pi.
\]
故
\[
\boxed{\int_0^1dy\int_{\pi-\arcsin y}^{\pi}f(x,y)\,dx}.
\]
\end{enumerate}
\examnote{含三角函数的换序要先看反三角函数所在象限；本题不能把下限误写成 \(\arcsin y\)。}
\end{solutionblock}

\begin{problemblock}
\textbf{2.} 将极坐标形式的累次积分改写为直角坐标积分。
\begin{enumerate}
\item \(\displaystyle \int_{\pi/4}^{\pi/2}d\theta\int_0^{2\sin\theta} f(r\cos\theta,r\sin\theta)\,r\,dr\).
\item \(\displaystyle \int_0^{\pi/4}d\theta\int_0^{2\cos\theta} f(r\cos\theta,r\sin\theta)\,r\,dr\).
\end{enumerate}
\end{problemblock}

\begin{solutionblock}
\analysis{先把极坐标边界化成圆，再结合角度限制确定区域。}
\begin{enumerate}
\item \(r=2\sin\theta\) 化为
\[
x^2+(y-1)^2=1.
\]
又 \(\pi/4\le\theta\le\pi/2\)，故区域在圆内且 \(y\ge x,x\ge0\)。按 \(x\) 积分可写成
\[
\boxed{\int_0^1dx\int_x^{1+\sqrt{1-x^2}}f(x,y)\,dy}.
\]
选择题答案为 D。
\item \(r=2\cos\theta\) 化为
\[
(x-1)^2+y^2=1.
\]
又 \(0\le\theta\le\pi/4\)，故区域在圆内且 \(0\le y\le x\)。按 \(y\) 积分为
\[
\boxed{\int_0^1dy\int_y^{1+\sqrt{1-y^2}}f(x,y)\,dx}.
\]
选择题答案为 B。
\end{enumerate}
\end{solutionblock}

\begin{problemblock}
\textbf{3.} 设 \(f(x,y)\) 连续，求
\[
\int_0^{\pi/4}d\theta\int_0^1f(r\cos\theta,r\sin\theta)\,r\,dr
\]
对应的直角坐标形式。
\end{problemblock}

\begin{solutionblock}
\analysis{这是单位圆第一象限内 \(0\le\theta\le\pi/4\) 的扇形，即 \(0\le y\le x\)。}
区域为
\[
x^2+y^2\le1,\qquad 0\le y\le x.
\]
按 \(y\) 先固定，
\[
0\le y\le\frac{\sqrt2}{2},\qquad y\le x\le\sqrt{1-y^2}.
\]
故
\[
\boxed{\int_0^{\sqrt2/2}dy\int_y^{\sqrt{1-y^2}}f(x,y)\,dx}.
\]
选择题答案为 C。
\end{solutionblock}

\begin{problemblock}
\textbf{4.} 设 \(f(x,y)\) 连续，判断
\[
\int_0^1dy\int_{-\sqrt{1-y^2}}^{1-y}f(x,y)\,dx
\]
的极坐标形式。
\end{problemblock}

\begin{solutionblock}
\analysis{左边界是单位圆左半弧，右边界是直线 \(x+y=1\)。}
区域在上半平面。极坐标下直线为
\[
r(\cos\theta+\sin\theta)=1.
\]
当 \(0\le\theta\le\pi/2\) 时，边界先遇到直线；当 \(\pi/2\le\theta\le\pi\) 时，边界为单位圆。因此
\[
\boxed{
\int_0^{\pi/2}d\theta\int_0^{1/(\cos\theta+\sin\theta)}
 f(r\cos\theta,r\sin\theta)\,r\,dr
+\int_{\pi/2}^{\pi}d\theta\int_0^1
 f(r\cos\theta,r\sin\theta)\,r\,dr}.
\]
选择题答案为 D。
\end{solutionblock}

\begin{problemblock}
\textbf{5.} 设区域 \(D\) 由 \(y=\sin x,\ x=\pm\frac{\pi}{2},\ y=1\) 围成，求
\[
\iint_D(xy^5-1)\,dxdy.
\]
\end{problemblock}

\begin{solutionblock}
\analysis{区域关于 \(y\) 轴对称，\(xy^5\) 关于 \(x\) 为奇函数。}
因此
\[
\iint_D xy^5\,dxdy=0.
\]
区域面积为
\[
S_D=\int_{-\pi/2}^{\pi/2}(1-\sin x)\,dx=\pi.
\]
故
\[
\iint_D(xy^5-1)\,dxdy=-S_D=\boxed{-\pi}.
\]
选择题答案为 D。
\end{solutionblock}

\begin{problemblock}
\textbf{6.} 设 \(f(x,y)\) 连续，且
\[
f(x,y)=xy+\iint_D f(x,y)\,dxdy,
\]
其中 \(D\) 由 \(y=0,y=x^2,x=1\) 围成，求 \(f(x,y)\)。
\end{problemblock}

\begin{solutionblock}
\analysis{把积分项看成常数 \(C\)，再由两边积分确定 \(C\)。}
令
\[
C=\iint_D f(x,y)\,dxdy.
\]
则
\[
f(x,y)=xy+C.
\]
又
\[
D:\quad 0\le x\le1,\quad 0\le y\le x^2,
\]
所以
\[
C=\int_0^1dx\int_0^{x^2}(xy+C)\,dy
=\frac1{12}+\frac{C}{3}.
\]
解得
\[
C=\frac18.
\]
故
\[
\boxed{f(x,y)=xy+\frac18}.
\]
选择题答案为 C。
\end{solutionblock}

\begin{problemblock}
\textbf{7.} 设 \(0<a<1\)，区域 \(D\) 由坐标轴、直线 \(x+y=a\) 与 \(x+y=1\) 围成，比较
\[
I=\iint_D\sin^2(x+y)\,d\sigma,\quad
J=\iint_D\ln^3(x+y)\,d\sigma,\quad
K=\iint_D(x+y)\,d\sigma.
\]
\end{problemblock}

\begin{solutionblock}
\analysis{令 \(u=x+y\)。在本区域内 \(a\le u\le1\)，横截长度为 \(u\)。}
因为 \(0<u<1\) 时
\[
\ln^3u<0,\qquad 0<\sin^2u<u,
\]
所以逐点比较得
\[
\ln^3(x+y)<\sin^2(x+y)<x+y.
\]
积分保持不等号，故
\[
\boxed{J<I<K}.
\]
选择题答案为 D。
\end{solutionblock}

\begin{problemblock}
\textbf{8.} 比较
\[
I=\iint_{|x|+|y|\le1}(x^2+y^3)\,d\sigma,\quad
J=\iint_{x^2+y^2\le1}(x^4-y^4)\,d\sigma,\quad
K=\iint_{x^2+y^2\le1}(x^3-y^2)\,d\sigma.
\]
\end{problemblock}

\begin{solutionblock}
\analysis{对称性是本题最快方法。}
在菱形区域中，\(y^3\) 为关于 \(y\) 的奇函数，积分为零，故
\[
I=\iint_{|x|+|y|\le1}x^2\,d\sigma>0.
\]
单位圆关于 \(x,y\) 对称，故
\[
J=\iint(x^4-y^4)\,d\sigma=0.
\]
又 \(x^3\) 关于 \(x\) 为奇函数，故
\[
K=-\iint_{x^2+y^2\le1}y^2\,d\sigma<0.
\]
于是
\[
\boxed{K<J<I}.
\]
选择题答案为 D。
\end{solutionblock}

\begin{problemblock}
\textbf{9.} 设
\[
I_1=\iint_D\frac{x+y}{4}\,d\sigma,\quad
I_2=\iint_D\sqrt{\frac{x+y}{4}}\,d\sigma,\quad
I_3=\iint_D\sqrt[3]{\frac{x+y}{4}}\,d\sigma,
\]
其中 \(D:(x-1)^2+(y-1)^2\le2\)。比较三者大小。
\end{problemblock}

\begin{solutionblock}
\analysis{圆盘上 \(0\le (x+y)/4\le1\)，比较 \(u,u^{1/2},u^{1/3}\) 即可。}
令
\[
u=\frac{x+y}{4}.
\]
在 \(D\) 上，\(x+y\) 的最小值为 \(0\)，最大值为 \(4\)，故 \(0\le u\le1\)。当 \(0<u<1\) 时
\[
u<u^{1/2}<u^{1/3}.
\]
因此
\[
\boxed{I_1<I_2<I_3}.
\]
选择题答案为 A。
\end{solutionblock}

\begin{problemblock}
\textbf{10.} 正方形 \(|x|\le1,|y|\le1\) 被对角线划分成四个区域 \(D_k\)，
\[
I_k=\iint_{D_k}y\cos x\,dxdy.
\]
求 \(\max I_k\)。
\end{problemblock}

\begin{solutionblock}
\analysis{在 \([-1,1]\) 上 \(\cos x>0\)，积分大小主要由 \(y\) 的符号决定。}
图中 \(D_1\) 位于上方，\(y>0\)，故 \(I_1>0\)。\(D_3\) 位于下方，故 \(I_3<0\)。左右两个区域关于 \(x\) 轴对称，\(y\cos x\) 关于 \(y\) 为奇函数，故
\[
I_2=I_4=0.
\]
所以最大者为
\[
\boxed{I_1}.
\]
选择题答案为 A。
\end{solutionblock}

\begin{problemblock}
\textbf{11.} \(D_k\) 是单位圆域在第 \(k\) 象限的部分，
\[
I_k=\iint_{D_k}(y-x)\,dxdy.
\]
判断哪个 \(I_k>0\)。
\end{problemblock}

\begin{solutionblock}
\analysis{第一象限中 \(x,y\) 对称，第二象限中 \(y-x\) 恒正。}
第一象限
\[
\iint_{D_1}y\,d\sigma=\iint_{D_1}x\,d\sigma,
\]
故 \(I_1=0\)。第二象限有 \(y\ge0,x\le0\)，且除边界外 \(y-x>0\)，故
\[
I_2>0.
\]
第三象限 \(y-x\) 的积分仍因对称为 \(0\)，第四象限 \(y-x<0\)。故答案为
\[
\boxed{I_2>0}.
\]
选择题答案为 B。
\end{solutionblock}

\begin{problemblock}
\textbf{12.} 已知
\[
\lim_{t\to0^+}\frac{\displaystyle\int_0^t dx\int_t^x e^{-y^2}\,dy}{t^\alpha}=\beta\ne0,
\]
求 \(\alpha,\beta\)。
\end{problemblock}

\begin{solutionblock}
\analysis{小区域极限先用 \(e^{-y^2}\sim1\) 抓主部。注意内层上限 \(x<t\)，积分为负。}
当 \(t\to0^+\) 时，
\[
\int_t^x e^{-y^2}\,dy\sim x-t.
\]
故分子
\[
\int_0^t(x-t)\,dx=-\frac{t^2}{2}.
\]
要使极限为非零常数，应取
\[
\alpha=2,\qquad \beta=-\frac12.
\]
选择题答案为 C。
\end{solutionblock}

\begin{problemblock}
\textbf{13.} 交换积分次序：
\[
\int_0^4dx\int_{\sqrt{4x-x^2}}^{2\sqrt{x}}f(x,y)\,dy.
\]
\end{problemblock}

\begin{solutionblock}
\analysis{边界 \(y=\sqrt{4x-x^2}\) 是圆 \((x-2)^2+y^2=4\) 的上半部分，\(y=2\sqrt{x}\) 是抛物线 \(x=y^2/4\)。水平切片时区域分成两块。}
原区域满足
\[
0\le x\le4,\qquad \sqrt{4x-x^2}\le y\le2\sqrt{x}.
\]
等价于
\[
y^2\le4x,\qquad y^2\ge4x-x^2.
\]
当 \(0\le y\le2\) 时，圆的不等式给出两段；当 \(2\le y\le4\) 时只受抛物线限制。故
\[
\boxed{
\int_0^2dy\left(
\int_{y^2/4}^{\,2-\sqrt{4-y^2}} f(x,y)\,dx
+\int_{2+\sqrt{4-y^2}}^{4} f(x,y)\,dx\right)
+\int_2^4dy\int_{y^2/4}^{4}f(x,y)\,dx }.
\]
\examnote{本题最易漏掉 \(0\le y\le2\) 时右侧的小区域。}
\end{solutionblock}

\begin{problemblock}
\textbf{14.} 交换积分次序：
\[
\int_0^2dx\int_x^{\sqrt{2x-x^2}}f(x,y)\,dy.
\]
\end{problemblock}

\begin{solutionblock}
\analysis{题面外层到 \(2\)，但当 \(1<x<2\) 时上限 \(\sqrt{2x-x^2}\) 小于下限 \(x\)，所以这是有向积分。应拆成正向区域减去反向区域。}
先写成
\[
I=\int_0^1dx\int_x^{\sqrt{2x-x^2}}f\,dy
-\int_1^2dx\int_{\sqrt{2x-x^2}}^x f\,dy.
\]
第一块区域为
\[
0\le y\le1,\qquad 1-\sqrt{1-y^2}\le x\le y.
\]
第二块区域分为
\[
0\le y\le1,\quad 1+\sqrt{1-y^2}\le x\le2,
\]
以及
\[
1\le y\le2,\quad y\le x\le2.
\]
因此
\[
\boxed{
I=\int_0^1dy\int_{1-\sqrt{1-y^2}}^y f(x,y)\,dx
-\int_0^1dy\int_{1+\sqrt{1-y^2}}^2f(x,y)\,dx
-\int_1^2dy\int_y^2f(x,y)\,dx }.
\]
\end{solutionblock}

\begin{problemblock}
\textbf{15.} 计算
\[
\int_0^1dx\int_{x^2}^{1}\frac{xy}{\sqrt{1+y^3}}\,dy.
\]
\end{problemblock}

\begin{solutionblock}
\analysis{先换序，内层对 \(x\) 积分即可。}
区域为
\[
0\le x\le1,\qquad x^2\le y\le1,
\]
即
\[
0\le y\le1,\qquad 0\le x\le\sqrt y.
\]
于是
\[
\begin{aligned}
I&=\int_0^1dy\int_0^{\sqrt y}\frac{xy}{\sqrt{1+y^3}}\,dx\\
&=\frac12\int_0^1\frac{y^2}{\sqrt{1+y^3}}\,dy
=\frac13(\sqrt2-1).
\end{aligned}
\]
故
\[
\boxed{\frac{\sqrt2-1}{3}}.
\]
\end{solutionblock}

\begin{problemblock}
\textbf{16.} 计算
\[
\int_0^1dy\int_{y/2}^{y}\cos x^2\,dx
+\int_1^2dy\int_{y/2}^{1}\cos x^2\,dx.
\]
\end{problemblock}

\begin{solutionblock}
\analysis{被积函数只含 \(x\)，换序后内层长度会变成 \(x\)。}
两个区域合并为
\[
0\le x\le1,\qquad x\le y\le2x.
\]
故
\[
I=\int_0^1dx\int_x^{2x}\cos x^2\,dy
=\int_0^1x\cos x^2\,dx
=\boxed{\frac12\sin1}.
\]
\end{solutionblock}

\begin{problemblock}
\textbf{17.} 计算
\[
\int_0^1dy\int_y^1(x^2-y^2)\,dx.
\]
\end{problemblock}

\begin{solutionblock}
\analysis{本题直接积分即可，也可先换序。}
\[
\begin{aligned}
I&=\int_0^1\left[\frac{x^3}{3}-y^2x\right]_{x=y}^{1}dy\\
&=\int_0^1\left(\frac13-y^2+\frac23y^3\right)dy
=\boxed{\frac16}.
\end{aligned}
\]
\end{solutionblock}

\begin{problemblock}
\textbf{18.} 计算
\[
\iint_{x^2+y^2\le1}\left((x+1)^2+2y^2\right)\,dxdy.
\]
\end{problemblock}

\begin{solutionblock}
\analysis{单位圆上奇函数项积分为零，且 \(\iint x^2=\iint y^2=\pi/4\)。}
展开得
\[
(x+1)^2+2y^2=x^2+2x+1+2y^2.
\]
故
\[
I=\frac{\pi}{4}+0+\pi+2\cdot\frac{\pi}{4}
=\boxed{\frac{7\pi}{4}}.
\]
\end{solutionblock}

\begin{problemblock}
\textbf{19.} 设 \(D=\{(x,y):0\le x\le1,0\le y\le1\}\)，讨论
\[
\iint_D\frac{dxdy}{x^2+y^2}.
\]
\end{problemblock}

\begin{solutionblock}
\analysis{唯一问题在原点，这是二重反常积分。}
在第一象限靠近原点的小扇形内，用极坐标有
\[
\frac{1}{x^2+y^2}\,dxdy=\frac1{r^2}\cdot r\,drd\theta=\frac{dr}{r}d\theta.
\]
而
\[
\int_0^\varepsilon\frac{dr}{r}=+\infty.
\]
所以该反常积分
\[
\boxed{\text{发散，且趋于 }+\infty}.
\]
\end{solutionblock}

\begin{problemblock}
\textbf{20.} 计算
\[
I=\int_0^{\pi/2}d\theta\int_0^{2\cos\theta}
\left[(r\cos\theta-1)^3+r\sin\theta\right]r\,dr.
\]
\end{problemblock}

\begin{solutionblock}
\analysis{区域是圆 \((x-1)^2+y^2\le1\) 的上半部分。}
令 \(u=x-1\)，则 \((x-1)^3=u^3\)。区域关于 \(u=0\) 对称，所以
\[
\iint_Du^3\,d\sigma=0.
\]
剩下
\[
I=\iint_D y\,d\sigma.
\]
这是半径为 \(1\) 的上半圆对 \(x\) 轴的一阶矩：
\[
\int_{-1}^1du\int_0^{\sqrt{1-u^2}}y\,dy
=\frac12\int_{-1}^1(1-u^2)\,du
=\boxed{\frac23}.
\]
\end{solutionblock}

\begin{problemblock}
\textbf{21.} 求极限
\[
\lim_{t\to0^+}\frac1{\sin^2t}\int_0^t dx\int_x^t e^{-(x-y)^2}\,dy.
\]
\end{problemblock}

\begin{solutionblock}
\analysis{小三角形面积给出主部，\(e^{-(x-y)^2}\to1\)，\(\sin^2t\sim t^2\)。}
\[
\int_0^t dx\int_x^t e^{-(x-y)^2}\,dy
\sim \int_0^t(t-x)\,dx=\frac{t^2}{2}.
\]
因此
\[
\boxed{\frac12}.
\]
\end{solutionblock}

\begin{problemblock}
\textbf{22.} 设
\[
f(t)=\int_0^t dx\int_x^{\sqrt{x}}\frac{\sin y}{y}\,dy,
\]
求 \(f(t)\) 在 \([0,\pi]\) 上的最大值。
\end{problemblock}

\begin{solutionblock}
\analysis{先由变上限求导确定最大点，再换序计算最大值。}
由微积分基本定理，
\[
f'(t)=\int_t^{\sqrt t}\frac{\sin y}{y}\,dy.
\]
当 \(0<t<1\) 时，\(\sqrt t>t\)，且 \(\sin y/y>0\)，故 \(f'(t)>0\)；当 \(1<t\le\pi\) 时，\(\sqrt t<t\)，故 \(f'(t)<0\)。所以最大值在 \(t=1\) 取得。

计算
\[
f(1)=\int_0^1dx\int_x^{\sqrt x}\frac{\sin y}{y}\,dy.
\]
区域等价于
\[
0\le y\le1,\qquad y^2\le x\le y.
\]
于是
\[
f(1)=\int_0^1(y-y^2)\frac{\sin y}{y}\,dy
=\int_0^1(1-y)\sin y\,dy
=\boxed{1-\sin1}.
\]
\end{solutionblock}

\begin{problemblock}
\textbf{23.} 求极限
\[
\lim_{n\to\infty}\frac1n\left(\int_{1/n}^1e^{-y^2}dy+\int_{2/n}^1e^{-y^2}dy+\cdots+\int_{(n-1)/n}^1e^{-y^2}dy\right).
\]
\end{problemblock}

\begin{solutionblock}
\analysis{这是函数 \(F(x)=\int_x^1e^{-y^2}dy\) 的 Riemann 和。}
\[
\lim_{n\to\infty}\frac1n\sum_{k=1}^{n-1}\int_{k/n}^1e^{-y^2}dy
=\int_0^1dx\int_x^1e^{-y^2}dy.
\]
换序得
\[
\int_0^1e^{-y^2}\left(\int_0^y dx\right)dy
=\int_0^1y e^{-y^2}dy
=\boxed{\frac{1-e^{-1}}2}.
\]
\end{solutionblock}

\begin{problemblock}
\textbf{24.} 求极限
\[
\lim_{t\to0^+}\frac1{t^6}\int_0^t dx\int_x^t\sin (xy)^2\,dy.
\]
\end{problemblock}

\begin{solutionblock}
\analysis{令 \(x=tu,y=tv\)，则 \((xy)^2=t^4u^2v^2\)，面积元给出 \(t^2\)。}
\[
\begin{aligned}
\int_0^t dx\int_x^t\sin (xy)^2\,dy
&=t^2\int_0^1du\int_u^1\sin(t^4u^2v^2)\,dv\\
&\sim t^6\int_0^1du\int_u^1u^2v^2\,dv.
\end{aligned}
\]
所以极限为
\[
\int_0^1u^2\frac{1-u^3}{3}\,du
=\frac13\left(\frac13-\frac16\right)
=\boxed{\frac1{18}}.
\]
\end{solutionblock}

\begin{problemblock}
\textbf{25.} 计算
\[
\int_{1/4}^{1/2}dy\int_{1/2}^{\sqrt y}e^{y/x}\,dx
+\int_{1/2}^{1}dy\int_y^{\sqrt y}e^{y/x}\,dx.
\]
\end{problemblock}

\begin{solutionblock}
\analysis{原区域由 \(x=\frac12,\ y=x^2,\ y=x\) 围成。换序后 \(y/x\) 适合直接积分。}
换序得
\[
I=\int_{1/2}^1dx\int_{x^2}^{x}e^{y/x}\,dy.
\]
令 \(u=y/x\)，则 \(dy=x\,du\)，内层上下限为 \(u=x\) 到 \(u=1\)。故
\[
\begin{aligned}
I&=\int_{1/2}^1x\int_x^1e^u\,du\,dx\\
&=\int_{1/2}^1x(e-e^x)\,dx\\
&=\frac{3e}{8}-\frac{\sqrt e}{2}.
\end{aligned}
\]
故
\[
\boxed{\frac{3e}{8}-\frac{\sqrt e}{2}}.
\]
\end{solutionblock}

\begin{problemblock}
\textbf{26.} 计算
\[
\iint_D|x^2+y^2-1|\,d\sigma,\qquad
D=\{(x,y):0\le x\le1,0\le y\le1\}.
\]
\end{problemblock}

\begin{solutionblock}
\analysis{单位正方形被第一象限单位圆弧分成两部分。}
在四分之一单位圆内，积分为 \(1-r^2\)；在其外，积分为 \(r^2-1\)。于是
\[
\begin{aligned}
I&=\int_0^{\pi/2}\int_0^1(1-r^2)r\,drd\theta
+\iint_{[0,1]^2}(x^2+y^2-1)\,d\sigma\\
&\quad-\int_0^{\pi/2}\int_0^1(r^2-1)r\,drd\theta\\
&=\frac{\pi}{8}+\left(\frac23-1\right)+\frac{\pi}{8}\\
&=\boxed{\frac{\pi}{4}-\frac13}.
\end{aligned}
\]
\end{solutionblock}

\begin{problemblock}
\textbf{27.} 计算
\[
\iint_D\max\{xy,1\}\,dxdy,\qquad
D=\{(x,y):0\le x\le2,0\le y\le2\}.
\]
\end{problemblock}

\begin{solutionblock}
\analysis{按双曲线 \(xy=1\) 分区。}
当 \(0\le x\le\frac12\) 时，\(xy\le1\) 恒成立；当 \(\frac12\le x\le2\) 时，以 \(y=1/x\) 分割。因此
\[
\begin{aligned}
I&=\int_0^{1/2}2\,dx
+\int_{1/2}^2\left(\int_0^{1/x}1\,dy+\int_{1/x}^2xy\,dy\right)dx\\
&=1+\int_{1/2}^2\left(2x+\frac1{2x}\right)dx\\
&=\boxed{\frac{19}{4}+\ln2}.
\end{aligned}
\]
\end{solutionblock}

\begin{problemblock}
\textbf{28.} 设
\[
D=\{(x,y):x^2+y^2\le2,\ x\ge0,\ y\ge0\},
\]
\([1+x^2+y^2]\) 表示不超过 \(1+x^2+y^2\) 的最大整数，计算
\[
\iint_Dxy[1+x^2+y^2]\,dxdy.
\]
\end{problemblock}

\begin{solutionblock}
\analysis{极坐标下取整项只依赖 \(r\)：\(0\le r<1\) 时为 \(1\)，\(1\le r<\sqrt2\) 时为 \(2\)。}
令 \(x=r\cos\theta,y=r\sin\theta\)，则
\[
0\le\theta\le\frac{\pi}{2},\quad 0\le r\le\sqrt2.
\]
所求积分为
\[
\int_0^{\pi/2}\cos\theta\sin\theta\,d\theta
\left(\int_0^1r^3\,dr+2\int_1^{\sqrt2}r^3\,dr\right).
\]
故
\[
I=\frac12\left(\frac14+\frac32\right)=\boxed{\frac78}.
\]
\end{solutionblock}

\begin{problemblock}
\textbf{29.} 计算
\[
\iint_D(x-y)\,dxdy,
\]
其中
\[
D=\{(x,y):(x-1)^2+(y-1)^2\le2,\ y\ge x\}.
\]
\end{problemblock}

\begin{solutionblock}
\analysis{绕圆心旋转坐标，使 \(y-x\) 成为一个坐标轴。}
令
\[
u=\frac{x+y}{\sqrt2},\qquad v=\frac{y-x}{\sqrt2}.
\]
则 \(x-y=-\sqrt2v\)，圆域变成
\[
(u-\sqrt2)^2+v^2\le2,
\]
且 \(y\ge x\) 即 \(v\ge0\)。因此
\[
I=-\sqrt2\iint_{\substack{(u-\sqrt2)^2+v^2\le2\\v\ge0}}v\,dudv.
\]
半径 \(R=\sqrt2\) 的上半圆对直径的一阶矩为 \(2R^3/3=4\sqrt2/3\)，故
\[
I=-\sqrt2\cdot\frac{4\sqrt2}{3}
=\boxed{-\frac83}.
\]
\end{solutionblock}

\begin{problemblock}
\textbf{30.} 计算
\[
I=\int_0^{\pi/4}\int_0^{\sec\theta}
r^2\sin\theta\sqrt{1-r^2\cos^2\theta}\,drd\theta.
\]
\end{problemblock}

\begin{solutionblock}
\analysis{虽然题面用极坐标给出，但化回直角坐标后更简单。}
令 \(x=r\cos\theta,y=r\sin\theta\)。区域为
\[
0\le x\le1,\qquad 0\le y\le x.
\]
又
\[
r^2\sin\theta\,drd\theta
=y\,dxdy,
\qquad
\sqrt{1-r^2\cos^2\theta}=\sqrt{1-x^2}.
\]
故
\[
I=\int_0^1dx\int_0^xy\sqrt{1-x^2}\,dy
=\frac12\int_0^1x^2\sqrt{1-x^2}\,dx.
\]
令 \(x=\sin t\)，得
\[
I=\frac12\int_0^{\pi/2}\sin^2t\cos^2t\,dt
=\boxed{\frac{\pi}{32}}.
\]
\end{solutionblock}

\begin{problemblock}
\textbf{31.} 计算
\[
\iint_D\frac{\sqrt{x^2+y^2}}{\sqrt{4a^2-x^2-y^2}}\,d\sigma,
\]
其中 \(D\) 是由曲线 \(y=-a+\sqrt{a^2-x^2}\ (a>0)\) 和直线 \(y=-x\) 围成的区域。
\end{problemblock}

\begin{solutionblock}
\analysis{曲线是圆 \(x^2+(y+a)^2=a^2\) 的上半部分，与直线 \(y=-x\) 围成第四象限内的小区域。}
极坐标下圆为
\[
r=-2a\sin\theta,
\]
区域为
\[
-\frac{\pi}{4}\le\theta\le0,\qquad 0\le r\le-2a\sin\theta.
\]
积分化为
\[
I=\int_{-\pi/4}^0d\theta\int_0^{-2a\sin\theta}
\frac{r^2}{\sqrt{4a^2-r^2}}\,dr.
\]
利用
\[
\int\frac{r^2}{\sqrt{R^2-r^2}}\,dr
=\frac{R^2}{2}\arcsin\frac rR-\frac r2\sqrt{R^2-r^2},
\]
取 \(R=2a\)，代入上限 \(r=-2a\sin\theta\)，得内层积分
\[
-2a^2\theta+2a^2\sin\theta\cos\theta.
\]
故
\[
I=a^2\left[-\theta^2+\sin^2\theta\right]_{-\pi/4}^0
=\boxed{a^2\left(\frac{\pi^2}{16}-\frac12\right)}.
\]
\end{solutionblock}

\begin{problemblock}
\textbf{32.} 计算
\[
\iint_D(x+y)^3\,dxdy,
\]
其中 \(D\) 由曲线 \(x=\sqrt{1+y^2}\) 与直线 \(x+\sqrt2y=0,\ x-\sqrt2y=0\) 围成。
\end{problemblock}

\begin{solutionblock}
\analysis{区域关于 \(x\) 轴对称，展开 \((x+y)^3\) 后含奇次 \(y\) 的项积分为零。}
区域可写为
\[
-1\le y\le1,\qquad \sqrt2|y|\le x\le\sqrt{1+y^2}.
\]
展开
\[
(x+y)^3=x^3+3x^2y+3xy^2+y^3.
\]
含 \(y\) 奇次的两项积分为零，故
\[
I=\int_{-1}^1dy\int_{\sqrt2|y|}^{\sqrt{1+y^2}}(x^3+3xy^2)\,dx.
\]
利用偶对称，
\[
\begin{aligned}
I&=2\int_0^1\left[
\frac{x^4}{4}+\frac{3}{2}y^2x^2
\right]_{\sqrt2y}^{\sqrt{1+y^2}}dy\\
&=2\int_0^1\left(\frac14+2y^2-\frac94y^4\right)dy\\
&=\boxed{\frac{14}{15}}.
\end{aligned}
\]
\end{solutionblock}

\begin{problemblock}
\textbf{33.} 计算
\[
\iint_D(x+y^2)\,dxdy,
\]
其中
\[
D=\{(x,y):x^2+y^2\le2x+2y\}.
\]
\end{problemblock}

\begin{solutionblock}
\analysis{区域是以 \((1,1)\) 为圆心、半径 \(\sqrt2\) 的圆盘。}
令 \(u=x-1,v=y-1\)，则区域为 \(u^2+v^2\le2\)。有
\[
\iint_Dx\,d\sigma=1\cdot \pi(\sqrt2)^2=2\pi.
\]
又
\[
y^2=(v+1)^2=v^2+2v+1,
\]
其中 \(\iint 2v\,d\sigma=0\)，
\[
\iint_{u^2+v^2\le2}v^2\,dudv=\frac{\pi R^4}{4}=\pi.
\]
故
\[
\iint_Dy^2\,d\sigma=\pi+2\pi=3\pi.
\]
所以
\[
\boxed{5\pi}.
\]
\end{solutionblock}

\begin{problemblock}
\textbf{34.} 求
\[
\iint_D(\sqrt{x^2+y^2}+y)\,d\sigma,
\]
其中 \(D\) 是由圆 \(x^2+y^2=4\) 和 \((x+1)^2+y^2=1\) 所围成的阴影区域。
\end{problemblock}

\begin{solutionblock}
\analysis{由图可知 \(D\) 为大圆盘去掉左侧小圆盘。区域关于 \(x\) 轴对称，\(y\) 项积分为零。}
故只需计算
\[
\iint_D\sqrt{x^2+y^2}\,d\sigma.
\]
大圆盘贡献为
\[
\int_0^{2\pi}d\theta\int_0^2r^2\,dr=\frac{16\pi}{3}.
\]
小圆 \((x+1)^2+y^2\le1\) 在极坐标中为
\[
\frac{\pi}{2}\le\theta\le\frac{3\pi}{2},\qquad 0\le r\le-2\cos\theta.
\]
其贡献为
\[
\int_{\pi/2}^{3\pi/2}d\theta\int_0^{-2\cos\theta}r^2\,dr
=\frac{32}{9}.
\]
因此
\[
\boxed{\frac{16\pi}{3}-\frac{32}{9}}.
\]
\end{solutionblock}

\begin{problemblock}
\textbf{36.} 计算
\[
\int_{\pi/4}^{3\pi/4}d\theta\int_0^{2\sin\theta}
\left(\sin\theta+\cos\theta\sqrt{1+r^2\sin^2\theta}\right)r^2\,dr.
\]
\end{problemblock}

\begin{solutionblock}
\analysis{积分区域关于 \(y\) 轴对称，含 \(\cos\theta\) 的部分化为关于 \(x\) 的奇函数，积分为零。}
令 \(x=r\cos\theta,y=r\sin\theta\)。区域关于 \(y\) 轴对称。第二项
\[
r^2\cos\theta\sqrt{1+r^2\sin^2\theta}\,drd\theta
=x\sqrt{1+y^2}\,dxdy
\]
关于 \(x\) 为奇函数，积分为零。

因此只需算
\[
\int_{\pi/4}^{3\pi/4}\int_0^{2\sin\theta}r^2\sin\theta\,drd\theta
=\frac83\int_{\pi/4}^{3\pi/4}\sin^4\theta\,d\theta.
\]
又
\[
\int_{\pi/4}^{3\pi/4}\sin^4\theta\,d\theta=\frac{3\pi}{16}+\frac12,
\]
故
\[
\boxed{\frac{\pi}{2}+\frac43}.
\]
\end{solutionblock}

\begin{problemblock}
\textbf{37.} 计算
\[
\int_{-1}^{1}dx\int_{|x|}^{1+\sqrt{1-x^2}}
(x^3+1)\sqrt{x^2+y^2}\,dy.
\]
\end{problemblock}

\begin{solutionblock}
\analysis{区域关于 \(y\) 轴对称，\(x^3\sqrt{x^2+y^2}\) 关于 \(x\) 为奇函数，积分为零。}
剩下
\[
\iint_D\sqrt{x^2+y^2}\,d\sigma.
\]
极坐标下区域为
\[
\frac{\pi}{4}\le\theta\le\frac{3\pi}{4},\qquad 0\le r\le2\sin\theta.
\]
故
\[
I=\int_{\pi/4}^{3\pi/4}d\theta\int_0^{2\sin\theta}r^2\,dr
=\frac83\int_{\pi/4}^{3\pi/4}\sin^3\theta\,d\theta.
\]
计算
\[
\int_{\pi/4}^{3\pi/4}\sin^3\theta\,d\theta
=\int_{\pi/4}^{3\pi/4}\sin\theta(1-\cos^2\theta)\,d\theta
=\left[-\cos\theta+\frac{\cos^3\theta}{3}\right]_{\pi/4}^{3\pi/4}
=\frac{5\sqrt2}{6}.
\]
故
\[
\boxed{\frac{20\sqrt2}{9}}.
\]
\end{solutionblock}

\begin{problemblock}
\textbf{38.} 设
\[
f(t)=\int_0^t dx\int_x^t y^2e^{-y^2}\,dy.
\]
证明对一切 \(t\in(-\infty,+\infty)\)，有
\[
0\le f(t)<\frac12.
\]
\end{problemblock}

\begin{solutionblock}
\analysis{把 \(f(t)\) 看成变上限函数求导，可得到一元函数表达式。}
由 Leibniz 公式，
\[
f'(t)=t\cdot t^2e^{-t^2}=t^3e^{-t^2},\qquad f(0)=0.
\]
因此
\[
f(t)=\int_0^t s^3e^{-s^2}\,ds
=\frac12\left[1-(t^2+1)e^{-t^2}\right].
\]
令 \(u=t^2\ge0\)。因为
\[
(u+1)e^{-u}\le1,
\]
且当 \(u>0\) 时严格小于 \(1\)，所以
\[
f(t)\ge0.
\]
又 \((u+1)e^{-u}>0\)，故
\[
f(t)<\frac12.
\]
命题得证。
\end{solutionblock}

\begin{problemblock}
\textbf{39.} 设 \(D=\{(x,y):0\le x\le2,0\le y\le2\}\)。
\begin{enumerate}
\item 计算 \(b=\iint_D|xy-1|\,d\sigma\)。
\item 设 \(f(x,y)\) 在 \(D\) 上连续，且
\[
\iint_Df(x,y)\,d\sigma=0,\qquad
\iint_Dxyf(x,y)\,d\sigma=1.
\]
证明存在 \((\xi,\eta)\in D\)，使
\[
|f(\xi,\eta)|\ge\frac1b.
\]
\end{enumerate}
\end{problemblock}

\begin{solutionblock}
\analysis{第一问按双曲线 \(xy=1\) 分割；第二问用积分绝对值估计和连续函数最大值。}
\begin{enumerate}
\item 当 \(0\le x\le\frac12\) 时，\(xy\le1\)；当 \(\frac12\le x\le2\) 时，以 \(y=1/x\) 分割。于是
\[
\begin{aligned}
b&=\int_0^{1/2}\int_0^2(1-xy)\,dy\,dx\\
&\quad+\int_{1/2}^2\left(\int_0^{1/x}(1-xy)\,dy+\int_{1/x}^{2}(xy-1)\,dy\right)dx\\
&=\frac34+\int_{1/2}^2\left(2x-2+\frac1x\right)dx\\
&=\boxed{\frac32+2\ln2}.
\end{aligned}
\]
\item 由条件
\[
1=\iint_Dxyf\,d\sigma-\iint_Df\,d\sigma
=\iint_D(xy-1)f(x,y)\,d\sigma.
\]
所以
\[
1\le \iint_D|xy-1|\,|f(x,y)|\,d\sigma
\le \max_D|f|\cdot \iint_D|xy-1|\,d\sigma
=b\max_D|f|.
\]
因此
\[
\max_D|f|\ge\frac1b.
\]
由于 \(f\) 连续，\(|f|\) 在闭区域 \(D\) 上能取到最大值，故存在 \((\xi,\eta)\in D\)，使
\[
|f(\xi,\eta)|\ge\frac1b.
\]
\end{enumerate}
\end{solutionblock}

\begin{problemblock}
\textbf{40.} 设 \(f(x),g(x)\) 在 \([0,1]\) 上连续，且同时单调增，证明
\[
\int_0^1f(x)g(x)\,dx\ge
\left(\int_0^1f(x)\,dx\right)
\left(\int_0^1g(x)\,dx\right).
\]
\end{problemblock}

\begin{solutionblock}
\analysis{这是积分形式的 Chebyshev 不等式。考研中常用“双重积分对称化”证明。}
因为 \(f,g\) 同时单调增，所以对任意 \(x,y\in[0,1]\)，有
\[
(f(x)-f(y))(g(x)-g(y))\ge0.
\]
对单位正方形积分，得
\[
\int_0^1\int_0^1(f(x)-f(y))(g(x)-g(y))\,dxdy\ge0.
\]
展开左边：
\[
\begin{aligned}
&\int_0^1\int_0^1(f(x)-f(y))(g(x)-g(y))\,dxdy\\
&=2\int_0^1f(x)g(x)\,dx
-2\left(\int_0^1f(x)\,dx\right)\left(\int_0^1g(x)\,dx\right).
\end{aligned}
\]
于是
\[
\int_0^1f(x)g(x)\,dx\ge
\left(\int_0^1f(x)\,dx\right)
\left(\int_0^1g(x)\,dx\right).
\]
证毕。
\examnote{若二者一个增一个减，则不等号方向相反。}
\end{solutionblock}

\section{本章小结}
第六章可见题号为 \(1\)--\(34\)、\(36\)--\(40\)，原书页面中未见第 \(35\) 题。本章核心方法是画区域、换序、对称性、极坐标与积分估计；其中第 \(13\)、\(14\) 题要特别注意区域分裂和有向积分。
